\chapter{Introduction}\label{ch:introduction}
%----------------------------------------------------------------------
\section{Purpose}
%----------------------------------------------------------------------

This document describes the design of the \ssat{} that will be
developed for computing worst-case needs for stack space. Its intended
readership includes primarily software engineers at AdaCore, as well
as users wishing to assess the qualification of the tool.

%----------------------------------------------------------------------
\section{Scope}
%----------------------------------------------------------------------

Stack overflows are a major threat to many computer applications and
run-time recovery techniques are not always available or appropriate.

In the safety-critical area, every possible effort must be made to
prevent such overflows from occurring in the first place, which calls
for methods to anticipate the worst-case needs for stack space.

The product to be developed will anticipate the stack space
requirements of applications and allow safe prevention of
overflow conditions.

The \ssat{} will analyze the application stack space consumption patterns
and compute worst-case bounds prior to execution time.

%----------------------------------------------------------------------
\section{Glossary}
%----------------------------------------------------------------------

\subsection{Acronyms}

\begin{description}
\item[COTS]  Commercial Off-The-Shelf. 
\item[GCC]   GNU Compiler Collection.
\item[GDB]   GNU debugger.
% \item[GMGPL] GNAT Modified General Public License.
\item[GNAT]  GNU Ada Translator
\item[GPL]   General Public License.
\item[GPS]   GNAT Programming Studio.
\item[HIS]   High integrity system.
\item[IDE]   Integrated Development Environment.
\end{description}

\subsection{Definition of terms}

\begin{description}
\item[High integrity system.] A system that must be trusted to work
  dependably and may result in unacceptable loss or harm
  otherwise. This includes safety-critical systems and other critical
  systems where, for example, security or financial considerations may
  be paramount.
\item[Cross-compilation system.] A compiler that builds programs for a
  machine different than the one the compiler is run on.
\item[Host platform.] The machine the toolchain programs will run on.
\item[Target platform.] The machine the toolchain builds programs for.
\end{description}

%----------------------------------------------------------------------
\section{References}
%----------------------------------------------------------------------

\begin{enumerate}[{R}1.]
\item \label{ALRM}
   ANSI/ISO/IEC-8652:1995. \textit{Ada~95 Reference Manual.}
   See also: S.T. Taft, R.A. Duff, E. Ploederer, \textit{Consolidated
     Ada Reference Manual}. LNCS 2219, Springer-Verlag.
\item  \label{HIS}
   ISO/IEC/TR-15942:2000 \textit{Guide for the use of the Ada programming
   language in high integrity systems.} 
\item \label{GNATUG}
      AdaCore. \textit{GNAT Pro User's Guide}. Version 5.04a. 2006.
\item \label{GNATRM}
      AdaCore. \textit{GNAT Pro Reference Manual}. Version 5.04a. 2006.
\item \label{GPS}
      AdaCore. \textit{GPS User's Guide}. Version 3.1.3. 2006.
\item \label{GDB}
      Free Software Foundation. \textit{Debugging with GDB}. Version 6.3 for GNAT Pro 5.04a. 2006.
\item \label{AdaQS}
  C. Ausnit-Hood, K.A. Johnson, R.G. Petit IV, and S.B. Opdahl.
  \textit{Ada 95 Quality and Style}.
  Springer-Verlag, LNCS 1344, 1995.
\item \label{GCC}
      Free Software Foundation. \textit{GCC, the GNU Compiler Collection}, \url{http://gcc.gnu.org}.
\item \label{VCG}
      \textit{Visualization of Compiler Graphs}, \url{http://rw4.cs.uni-sb.de/users/sander/html/gsvcg1.html}.
\item \label{GCCSTACK}
      Eric Botcazou, Cyrille Comar, and Olivier Hainque, \textit{Compile-time stack requirements analysis with GCC}. In Proceedings of the 3rd GCC developers summit. 2005.
\end{enumerate}

%----------------------------------------------------------------------
\section{Document overview}
%----------------------------------------------------------------------

This document is organized as follows:
chapter \ref{ch:description} contains a general description of the
tool. Chapter \ref{ch:overview} provides a brief overview of the
different software components that will make up \ssat{}. These
software components are defined in more detail in chapter
\ref{ch:toplevel}. Finally, the detailed implementation of the
different run-time packages can be found in chapter \ref{ch:detailed}.

% LocalWords:  AdaCore GCC GDB GMGPL GPL GPS toolchain Ploederer LNCS Ausnit
% LocalWords:  Petit Opdahl Botcazou Cyrille Comar Hainque
